\chapter{Bermudéennes et remboursables par anticipation}\label{ch:rc:bermudans-and-callables}

En 2014, le législateur taïwanais a exempté du plafond applicable aux investissements à l'étranger de ses assureurs-vie les obligations libellées en devises
cotées sur l'île, et les assureurs, qui disposaient de centaines de milliards de dollars à placer,
se sont mis à acheter des obligations zéro-coupon en dollars de trente ans
que les banques émettrices pouvaient rembourser chaque année après les toutes premières années.
Les investisseurs obtenaient un rendement nettement supérieur à la courbe des bons du Trésor
en rémunération de la vente de ce droit. Les émetteurs ne le gardaient pas~: ils échangeaient
chaque obligation contre un financement à taux variable auprès d'un dealer, et le swap portait le droit de remboursement, si bien que les dealers
se retrouvaient longs d'un large stock d'options bermudéennes sur des taux à longue échéance. Pour couvrir
cette position, ils vendaient de la volatilité sur le marché des swaptions, année après année. L'instrument central
est la \omterm{def:rc:bermudans-and-callables:bermudan}{swaption bermudéenne}~: une option permettant d'entrer dans un swap à l'une quelconque de plusieurs dates.
Elle vaut plus que la plus chère des swaptions européennes qu'elle peut devenir, d'un montant qu'aucune couverture européenne ne capture.
Ce chapitre la valorise sur l'arbre du chapitre~7,
par simulation de Monte Carlo à régression, mesure ce montant et
valorise l'obligation zéro-coupon remboursable par anticipation (callable) que les assureurs ont achetée.

\section{Swaptions bermudéennes et leurs swaptions européennes co-terminales}

\begin{definition}[Swaption bermudéenne]\label{def:rc:bermudans-and-callables:bermudan}
Une \emph{swaption bermudéenne}\index{swaption bermudéenne} confère le droit, à l'une quelconque
d'un ensemble de dates d'exercice $T_{e_1}<\dots<T_{e_m}$, d'entrer dans le swap qui débute à
cette date et se termine à une date finale fixe $T_n$, à un taux fixe $K$. Une
bermudéenne receveuse «~ten-non-call-one~» peut être exercée chaque année, de l'année un
à l'année dix, pour entrer dans un swap receveur se terminant à l'année dix. Il s'agit d'un droit d'exercice bermudéen au
sens du livre~5, chapitre~6, appliqué aux swaps.
\end{definition}

\begin{definition}[Swaption co-terminale]\label{def:rc:bermudans-and-callables:coterminal}
Les \emph{swaptions co-terminales}\index{swaption co-terminale} d'une bermudéenne sont les
swaptions européennes qui expirent à chacune de ses dates d'exercice, pour entrer dans le swap se terminant
à sa date finale~: les européennes qu'elle peut devenir.
\end{definition}

\begin{definition}[Option d'échange]\label{def:rc:bermudans-and-callables:switch}
L'\emph{option d'échange}\index{option d'échange} (switch option) d'une bermudéenne est sa valeur diminuée de la
valeur de sa swaption européenne co-terminale la plus chère~: la valeur de la faculté de
choisir la date d'exercice après avoir observé les taux.
\end{definition}

\begin{proposition}[Bornes issues des européennes]\label{prop:rc:bermudans-and-callables:bounds}
Une bermudéenne vaut au moins sa swaption européenne co-terminale la plus chère et au plus
la somme de toutes.
\end{proposition}
\begin{proof}
Le détenteur peut s'engager sur une seule date d'exercice, ce qui donne l'européenne~; toute
politique d'exercice n'exerce qu'à une date au plus, de sorte que le paiement est dominé par le
somme des paiements de toutes les européennes.
\end{proof}

La borne inférieure est celle qu'une couverture statique par des européennes garantit~; l'\omterm{def:rc:bermudans-and-callables:switch}{option d'échange}
dépend de la manière dont les taux des swaps co-terminaux évoluent ensemble, ce que seul un modèle
de toute la courbe permet de valoriser.

\section{Exercice par induction rétrograde sur un arbre}

\begin{method}[Une bermudéenne sur l'arbre de Hull--White]\label{met:rc:bermudans-and-callables:tree}
Construire l'arbre du chapitre~7 jusqu'à la date finale. Remonter depuis la dernière date d'exercice~:
à chaque date d'exercice, calculer en chaque nœud la valeur du swap sous-jacent
(ses obligations zéro-coupon remontées sur l'arbre), puis remplacer la valeur de continuation
par le maximum des deux~; entre les dates d'exercice, actualiser le long des
branches. La valeur à la racine est le prix de la bermudéenne~; les nœuds où
l'exercice l'emporte tracent la frontière d'exercice.
\end{method}

\omcode{code/firm/bermudan/firm_bermudan.py}{29}{41}{Une swaption bermudéenne par induction rétrograde sur l'arbre.}\label{lst:rc:bermudans-and-callables:tree}

\begin{example}[Une receveuse ten-non-call-one]\label{ex:rc:bermudans-and-callables:tenno}
Sur la courbe OIS en euros du chapitre~2, avec Hull--White à $\kappa=3~\%$ et $\sigma=86$ points
de base (calibration à une volatilité du chapitre~7), une bermudéenne receveuse au taux au pair à dix ans
de 2.585~\%, exerçable chaque année de l'année un à l'année neuf, vaut 358.4 points de base
de notionnel. Ses européennes co-terminales valent 209.5, 245.3, 249.5, 237.8,
214.1, 184.3, 147.8, 104.3 et 54.9 points de base (arbre, à 1.4 près de la formule
de Jamshidian~; la \cref{fig:rc:bermudans-and-callables:europeans})~; la plus chère est
celle de trois ans. L'\omterm{def:rc:bermudans-and-callables:switch}{option d'échange} vaut 109.0 points de base, soit 44~\% de la meilleure
européenne.
\end{example}

\begin{omfigure}
\begin{tikzpicture}
\begin{axis}[width=11.5cm,height=5.4cm,ybar,bar width=14pt,
  xlabel={date d'exercice (années)~; le swap se termine à l'année dix},ylabel={valeur (pb de notionnel)},
  tick label style={font=\small},label style={font=\small},
  xmin=0.4,xmax=9.6,ymin=0,ymax=420,xtick={1,2,3,4,5,6,7,8,9},grid=major,grid style={gray!25},
  table/col sep=comma]
\addplot[fill=omDef!55,draw=omDef] table[x=e,y=euro]{figdata/rates-credit-risk/09-bermudans-and-callables/europeans.csv};
\addplot[omThm,very thick,sharp plot,mark=none,forget plot] table[x=e,y=berm]{figdata/rates-credit-risk/09-bermudans-and-callables/europeans.csv};
\node[font=\footnotesize,text=omThm,above] at (axis cs:7,365) {Bermudéenne};
\end{axis}
\end{tikzpicture}

\omcaption{Les swaptions receveuses co-terminales d'une bermudéenne ten-non-call-one en euros
au taux au pair à dix ans, et la bermudéenne elle-même (la ligne). L'écart entre la
ligne et la barre la plus haute est l'\omterm{def:rc:bermudans-and-callables:switch}{option d'échange}. Données~: courbe du chapitre~2,
Hull--White avec $\kappa=3~\%$ et $\sigma=86$ points de base~; le tutoriel
du chapitre.}\label{fig:rc:bermudans-and-callables:europeans}
\end{omfigure}

La frontière d'exercice de la receveuse est un taux court en dessous duquel l'exercice
l'emporte sur l'attente (la \cref{fig:rc:bermudans-and-callables:boundary}). Au début, le taux
doit baisser fortement, jusqu'à 0.82~\% à l'année un, car l'exercice abandonne huit occasions
supplémentaires~; à l'année neuf, l'exercice vaut la peine en dessous de 2.30~\%, proche du strike.

\begin{omfigure}
\begin{tikzpicture}
\begin{axis}[width=11.5cm,height=5.2cm,
  xlabel={date d'exercice (années)},ylabel={taux (\%)},
  tick label style={font=\small},label style={font=\small},
  xmin=0.5,xmax=9.5,ymin=0.5,ymax=3.3,xtick={1,2,3,4,5,6,7,8,9},grid=major,grid style={gray!25},
  legend columns=3,legend style={font=\footnotesize,at={(0.5,-0.25)},anchor=north,draw=none,fill=none,/tikz/every even column/.append style={column sep=8pt}},
  table/col sep=comma]
\addplot[omThm,very thick,mark=*] table[x=e,y=rstar]{figdata/rates-credit-risk/09-bermudans-and-callables/boundary.csv};
\addplot[omDef,thick,dashed] table[x=e,y=strike]{figdata/rates-credit-risk/09-bermudans-and-callables/boundary.csv};
\addplot[omProp,thick,dotted] table[x=e,y=fwd]{figdata/rates-credit-risk/09-bermudans-and-callables/boundary.csv};
\legend{exercice en dessous de ce taux court,strike 2.585~\%,taux forward à un an}
\end{axis}
\end{tikzpicture}

\omcaption{Frontière d'exercice de la \omterm{def:rc:bermudans-and-callables:bermudan}{swaption bermudéenne} receveuse sur l'arbre~: le taux court le plus élevé
auquel exercer dans le swap restant est préférable à la poursuite. Elle monte
vers le strike à mesure que les dates d'exercice se raréfient. Données~: le tutoriel
du chapitre.}\label{fig:rc:bermudans-and-callables:boundary}
\end{omfigure}

\begin{example}[Retour à la moyenne et l'option d'échange]\label{ex:rc:bermudans-and-callables:kappa}
En recalibrant $\sigma$ de sorte que la swaption européenne cinq ans sur cinq ans conserve son prix, la
\omterm{def:rc:bermudans-and-callables:bermudan}{swaption bermudéenne} et son \omterm{def:rc:bermudans-and-callables:switch}{option d'échange} dépendent de $\kappa$ (tableau ci-dessous)~: un retour à la moyenne
plus fort décorrèle les taux des dates d'exercice successives et rend le choix
de la date plus précieux, tandis que la meilleure européenne bouge à peine.
\end{example}

\begin{center}
\small
\begin{tabular}{rrrrr}
\toprule
$\kappa$ & $\sigma$ (bp) & Bermudéenne (bp) & meilleure européenne (bp) & échange (bp) \\
\midrule
1~\% & 78.1 & 351.0 & 249.2 & 101.9 \\
3~\% & 86.0 & 358.4 & 249.5 & 109.0 \\
6~\% & 98.9 & 369.7 & 249.7 & 120.0 \\
\bottomrule
\end{tabular}

\footnotesize Retour à la moyenne, volatilité recalibrée et receveuse à dix ans non remboursable à un an.
Données~: le tutoriel du chapitre.
\end{center}

\section{Exercice par régression dans Monte Carlo}

Les arbres ne servent qu'à un ou deux facteurs~; un modèle de marché (chapitre~8) ou tout modèle
à nombreuses variables d'état se simule, et la simulation avance dans le temps tandis que l'exercice
se décide à rebours. La méthode de Longstaff--Schwartz (livre~5, chapitre~23)
estime la valeur de continuation en régressant, à chaque date d'exercice, les
flux futurs actualisés des trajectoires dans la monnaie sur des fonctions de l'état.

\begin{method}[Exercice par régression pour une swaption bermudéenne]\label{met:rc:bermudans-and-callables:lsm}
Simuler l'état sous une mesure commode (ici la probabilité forward terminale
de l'année dix, numéraire $P(t,T_n)$) aux dates d'exercice. En remontant depuis la dernière date~:
calculer la valeur d'exercice en unités de numéraire~; sur les trajectoires dans la monnaie, régresser
le flux futur réalisé sur $1$, la valeur d'exercice et son carré~; exercer
lorsque la valeur d'exercice dépasse la continuation ajustée. Ajuster les régressions sur un
ensemble de trajectoires et appliquer la règle à un ensemble indépendant~: le résultat est le prix d'une
politique réalisable, donc une borne inférieure.
\end{method}

\omcode{code/firm/bermudan/firm_bermudan.py}{121}{142}{L'étape de régression de Longstaff--Schwartz, exécutée à rebours à travers les dates d'exercice.}\label{lst:rc:bermudans-and-callables:lsm}

\begin{example}[L'arbre et la régression concordent]\label{ex:rc:bermudans-and-callables:lsm}
Avec 40\,000 trajectoires sous la probabilité forward de dix ans (antithétiques, pas demi-mensuels)
le prix par régression de la receveuse à dix ans non remboursable à un an est de 358.1 points de base avec une
erreur type de 1.9, contre 358.4 sur l'arbre.
\end{example}

\section{Bornes supérieures et option d'échange}

Un prix par régression est une borne inférieure~; l'écart qui le sépare de la vraie valeur
se mesure par une borne supérieure duale (livre~5, chapitre~23), qui retranche une
martingale du payoff et prend le maximum sur les dates le long de chaque trajectoire~:
l'écart entre les deux bornes certifie la politique, et son calcul (par simulation imbriquée)
relève de la validation du pricer. Une fois l'écart petit, les
risques pratiques se situent ailleurs~: dans la structure de corrélation et de volatilité
du modèle, qui fixe la taille de l'\omterm{def:rc:bermudans-and-callables:switch}{option d'échange}.

\begin{remark}[Couvrir une bermudéenne]\label{rem:rc:bermudans-and-callables:hedge}
Le vega d'une bermudéenne repose sur ses \omterm{def:rc:bermudans-and-callables:coterminal}{swaptions co-terminales}, réparties à peu près
en proportion de la probabilité d'exercer dans chacune~; son \omterm{def:rc:bermudans-and-callables:switch}{option d'échange} est un pari sur la
décorrélation des taux de swap successifs que aucune européenne ne capture. Les desks couvrent le
vega et le carry des européennes et conservent en réserve le reste~; un changement de modèle (le $\sigma(t)$
du chapitre~7, le $\beta$ d'un modèle de marché) modifie la valeur de l'\omterm{def:rc:bermudans-and-callables:switch}{option d'échange}.
\end{remark}

\section{L'activité d'émission d'obligations remboursables par anticipation}

\begin{definition}[Obligation Formosa]\label{def:rc:bermudans-and-callables:formosa}
Une \emph{obligation Formosa}\index{obligation Formosa} est une obligation libellée dans une devise
étrangère et émise et cotée à Taïwan. À partir de 2014, la plupart des Formosa en dollars étaient
à longue échéance (vingt ou trente ans) et remboursables par anticipation à l'initiative de l'émetteur, beaucoup avec des
coupons zéro, vendues à des assureurs-vie taïwanais.
\end{definition}

\begin{definition}[Range accrual remboursable par anticipation]\label{def:rc:bermudans-and-callables:rangeaccrual}
Un \emph{range accrual remboursable par anticipation}\index{range accrual remboursable par anticipation} est un titre
remboursable par anticipation dont le coupon ne s'accumule que les jours où un taux de référence reste dans une fourchette~; le
droit de remboursement de l'émetteur en fait une bermudéenne sur un portefeuille d'options digitales, et il est
évalué avec les mêmes méthodes d'exercice, sur un modèle qui ajuste aussi le smile.
\end{definition}

\begin{omfigure}
\begin{tikzpicture}[font=\small,
  box/.style={draw,thick,rounded corners=2pt,align=center,minimum height=1.1cm,minimum width=2.8cm,inner sep=4pt},
  arr/.style={-{Stealth[length=2.5mm]},thick}]
\node[box,draw=omDef,fill=omDef!8] (inv) at (0,0) {assureur-vie\\(investisseur)};
\node[box,draw=omThm,fill=omThm!8] (iss) at (4.8,0) {banque émettrice};
\node[box,draw=omProp,fill=omProp!8] (dlr) at (9.6,0) {desk de taux\\du dealer};
\node[box,draw=gray,fill=gray!8] (mkt) at (9.6,-2.6) {marché des swaptions};
\draw[arr] (inv.15) -- node[above=4pt,align=center] {cash à l'émission} (iss.165);
\draw[arr] (iss.195) -- node[below=4pt,align=center] {coupon zéro, remboursable} (inv.345);
\draw[arr] (iss.15) -- node[above=4pt,align=center] {swap + droit de remboursement} (dlr.165);
\draw[arr] (dlr.195) -- node[below=4pt,align=center] {financement à taux variable} (iss.345);
\draw[arr,omExo] (dlr) -- node[right=4pt,align=left] {vend de la volatilité\\pour couvrir la bermudéenne} (mkt);
\end{tikzpicture}

\omcaption{Les flux d'une émission Formosa remboursable. L'investisseur vend le droit de
remboursement à l'émetteur contre un rendement plus élevé~; l'émetteur le transmet à un dealer dans le
swap qui transforme l'obligation en financement à taux variable~; le dealer, long d'une bermudéenne,
vend de la volatilité pour la couvrir.}\label{fig:rc:bermudans-and-callables:flows}
\end{omfigure}

\begin{dated}{2026-09}{Le marché Formosa}\label{dat:rc:bermudans-and-callables:formosa}
En mai 2014, le parlement de Taïwan a exclu les obligations en devises émises localement
du plafond de 45~\% sur les placements à l'étranger des assureurs, débloquant des fonds déclarés à
586~milliards de dollars~; les billets en dollars vendus dans les mois suivants étaient tous remboursables par anticipation,
d'échéances de vingt ou trente ans, et 74~\% ne versaient aucun coupon (données
Bloomberg). La presse spécialisée décrit les flux de couverture qui en résultent comme un
apport régulier de volatilité en dollars de long terme, qui a contribué à la maintenir basse.
\end{dated}

\section{Tutoriel~: une bermudéenne sous trois angles}\label{tut:rc:bermudans-and-callables:tutorial}

\begin{tutorial}
\textbf{Objectif.} Évaluer une swaption receveuse ten-non-call-one sur l'arbre et par régression,
mesurer son \omterm{def:rc:bermudans-and-callables:switch}{option d'échange} et sa frontière d'exercice. \textbf{État final~:}
les \cref{fig:rc:bermudans-and-callables:europeans,fig:rc:bermudans-and-callables:boundary}
et les chiffres des \cref{ex:rc:bermudans-and-callables:tenno,ex:rc:bermudans-and-callables:lsm}.
\begin{tutsteps}
\item \textbf{Arbre}~: \texttt{bermudan\_summary()} construit l'arbre du chapitre~7 avec
des pas semi-mensuels et évalue la bermudéenne et ses neuf européennes.
\item \textbf{Frontière}~: \texttt{boundary()}.
\item \textbf{Régression}~: \texttt{lsm\_check()} sur 40\,000 trajectoires.
\item \textbf{Retour à la moyenne}~: \texttt{mean\_reversion\_table()}~;
\texttt{fig\_rc\_bermudan.py} écrit les graphiques.
\end{tutsteps}
\textbf{Pour aller plus loin.} Ajouter $x^2$ et le taux à un an à la
base de régression et voir si le prix augmente~; évaluer la bermudéenne sur le
modèle du chapitre~7 à $\sigma$ par morceaux, par Monte-Carlo.
\end{tutorial}

\section{Construction~: moteurs de valorisation des bermudéennes et des produits remboursables}\label{bld:rc:bermudans-and-callables:bermudan}

\begin{build}
\bfield{Rôle} Les produits remboursables par anticipation du livre~: \omterm{def:rc:bermudans-and-callables:bermudan}{swaptions bermudéennes}, obligations
remboursables, billets compris, et les décisions d'exercice des prêts immobiliers du chapitre~12.
\bfield{Interface} \texttt{bermudan\_tree(tree, exercises, final, K, receiver)}~;
\texttt{exercise\_boundary}~; \texttt{callable\_zero\_tree(tree, maturity,
accretion, calls)}~; \texttt{simulate\_x}, \texttt{bermudan\_lsm(model, exercises,
final, K, paths)}~; \texttt{switch\_option}.
\bfield{Règles} Le \texttt{firm\_shortrate} du chapitre~7 est importé~; la régression est ajustée et
appliquée sur des ensembles de trajectoires indépendants~; graines fixées~; numpy seul.
\bfield{Tests d'acceptation} \texttt{code/firm/bermudan/tests/}~: une seule date
d'exercice donne l'européenne~; la bermudéenne se situe entre la meilleure européenne et la
somme~; la régression concorde avec l'arbre~; un zéro-coupon qui capitalise bien au-dessus des taux est
remboursé dès la première date, et un zéro-coupon capitalisé à zéro presque jamais~; la
frontière d'exercice se situe sous le strike.
\bfield{Pour aller plus loin} Une borne supérieure duale par simulation imbriquée~; les bermudéennes dans le modèle de
marché du chapitre~8~; \omterm{def:rc:bermudans-and-callables:rangeaccrual}{range accruals remboursables par anticipation}.
\end{build}

\begin{omsources}
\item F.\ Longstaff and E.\ Schwartz, ``Valuing American options by simulation: a
simple least-squares approach'', \emph{Review of Financial Studies} 14(1), 2001.
\item L.\ Andersen and M.\ Broadie, ``Primal-dual simulation algorithm for pricing
multidimensional American options'', \emph{Management Science} 50(9), 2004.
\item Insurance Journal (Bloomberg News), ``Goldman leads Wall Street tapping
Taiwan cash pile'', 10 septembre 2014~; Risk.net, ``Formosas, the Fed, and the
billion-dollar Bermudan trade''.
\item L.\ Andersen and V.\ Piterbarg, \emph{Interest Rate Modeling}, volume 3,
Atlantic Financial Press, 2010, chapitre~19.
\end{omsources}

\section{Exercices}

\begin{exercise}[$\star$]\label{exo:rc:bermudans-and-callables:1}
En utilisant l'\cref{ex:rc:bermudans-and-callables:tenno}, donner les bornes inférieure et supérieure
de la \cref{prop:rc:bermudans-and-callables:bounds} et vérifier que la bermudéenne se situe
entre elles.
\end{exercise}

\begin{exercise}[$\star$]\label{exo:rc:bermudans-and-callables:2}
Pourquoi l'européenne à trois ans est-elle la plus chère, et non celle à un an
portant le swap le plus long~?
\end{exercise}

\begin{exercise}[$\star$]\label{exo:rc:bermudans-and-callables:3}
Un investisseur achète un zéro-coupon de trente ans remboursable par anticipation plutôt
qu'un zéro ordinaire. Qu'a-t-elle vendu, et à qui ce droit revient-il finalement~?
\end{exercise}

\begin{exercise}[$\star\star$]\label{exo:rc:bermudans-and-callables:4}
Lire la \cref{fig:rc:bermudans-and-callables:boundary}. Pourquoi la frontière monte-t-elle
vers le strike, et pourquoi reste-t-elle sous le taux forward à un an~?
\end{exercise}

\begin{exercise}[$\star\star$]\label{exo:rc:bermudans-and-callables:5}
Expliquer, à partir du tableau de l'\cref{ex:rc:bermudans-and-callables:kappa}, pourquoi
l'\omterm{def:rc:bermudans-and-callables:switch}{option d'échange} augmente avec $\kappa$ alors que la meilleure européenne varie à peine.
\end{exercise}

\begin{exercise}[$\star\star$]\label{exo:rc:bermudans-and-callables:6}
Pourquoi faut-il ajuster la régression sur un ensemble de trajectoires et l'appliquer à un autre pour
donner une borne inférieure~?
\end{exercise}

\begin{exercise}[$\star\star\star$]\label{exo:rc:bermudans-and-callables:7}
\emph{Programmation.} Évaluer le zéro-coupon remboursable thirty-non-call-five du problème du week-end
à son rendement à capitalisation au pair, et trouver le rendement du zéro-coupon classique sur la même
courbe.
\end{exercise}

\begin{exercise}[$\star\star\star$]\label{exo:rc:bermudans-and-callables:8}
\emph{Trouver la faille.} «~Nous couvrons notre bermudéenne par la \omterm{def:rc:bermudans-and-callables:coterminal}{swaption co-terminale} européenne la plus chère,
à notionnel égal, donc nous sommes couverts.~»
\end{exercise}

\section{Problème~: la couverture Formosa}

\begin{problem}[{Problème du week-end --- un zéro-coupon remboursable de trente ans}]\label{pb:rc:bermudans-and-callables:1}
Une banque émet une \omterm{def:rc:bermudans-and-callables:formosa}{obligation Formosa} zéro-coupon en dollars de trente ans, remboursable chaque année à partir de
la cinquième année à sa valeur capitalisée, et la vend au pair à un assureur-vie. Elle l'échange
avec un dealer, qui reprend le droit de remboursement. Évaluer l'obligation avec le \omterm{def:rc:short-rate-models:hw}{modèle de Hull--White} ($\kappa=3\%$,
$\sigma=90$ points de base) sur la courbe SOFR du chapitre~1 (convexe monotone), arbre mensuel.

\medskip\noindent\textbf{Partie I --- L'obligation.}
\begin{enumerate}
\item Donner le rendement zéro-coupon à trente ans de la courbe (capitalisation annuelle).
\item Donner le rendement de capitalisation pour lequel le zéro-coupon remboursable vaut le pair.
\item Donner la valeur du zéro-coupon simple à ce taux de capitalisation, par unité émise.
\item Donner la valeur du call détenu par l'émetteur.
\item Pourquoi l'investisseur accepte-t-il cela~?
\end{enumerate}

\medskip\noindent\textbf{Partie II --- L'option.}
\begin{enumerate}[resume]
\item Quelle date de remboursement unique vaut le plus en tant qu'option européenne, et pour quel montant~?
\item Donner la prime de l'option bermudéenne par rapport à cette européenne.
\item Pourquoi la première date de remboursement est-elle la plus valorisée parmi les européennes ici~?
\item Que détient le dealer en position longue, et contre quoi~?
\item Pourquoi le dealer ne garde-t-il pas simplement l'option~?
\end{enumerate}

\medskip\noindent\textbf{Partie III --- La couverture.}
\begin{enumerate}[resume]
\item Donner le vega du call par point de base de volatilité, par tranche de 1 milliard de dollars~USD.
\item Quelles swaptions le dealer vend-il pour le couvrir~?
\item Si dix émissions de 1 milliard de dollars~USD de ce type arrivent dans un mois, quel vega entre sur le marché~?
\item Que se passe-t-il pour la volatilité des swaptions à longue échéance, et pour les valorisations des dealers~?
\item Lorsque les taux baissent et que les obligations sont remboursées, que devient la couverture~?
\end{enumerate}

\medskip\noindent\textbf{Partie IV --- Jugement.}
\begin{enumerate}[resume]
\item Quels risques l'assureur conserve-t-il~?
\item Pourquoi la transaction est-elle concentrée sur peu de dealers, et pourquoi cela importe-t-il~?
\item Quel risque de modèle le dealer porte-t-il~?
\item Énoncer le \emph{résultat nommé}~: le rendement de capitalisation au pair, la prime de l'option bermudéenne par rapport
à la meilleure européenne et le vega par tranche de 1 milliard de dollars~USD.
\item En une phrase~: qui vend de la volatilité dans cette chaîne, et pourquoi~?
\end{enumerate}
\end{problem}

\section{Questions d'entretien}

\begin{interviewq}[$\star$ \iqroles{trader, banque}]\label{iq:rc:bermudans-and-callables:1}
Qu'est-ce qu'une \omterm{def:rc:bermudans-and-callables:bermudan}{swaption bermudéenne}, et pourquoi vaut-elle plus que chacune de ses européennes~?
\end{interviewq}

\begin{interviewq}[$\star\star$ \iqroles{chercheur, développeur}]\label{iq:rc:bermudans-and-callables:2}
Expliquez la méthode de Longstaff--Schwartz pour une \omterm{def:rc:bermudans-and-callables:bermudan}{swaption bermudéenne}. Pourquoi s'agit-il d'une borne inférieure~?
\end{interviewq}

\begin{interviewq}[$\star\star$ \iqroles{chercheur}]\label{iq:rc:bermudans-and-callables:3}
Quel paramètre d'un modèle à un facteur gouverne l'\omterm{def:rc:bermudans-and-callables:switch}{option d'échange}, et comment~?
\end{interviewq}

\begin{interviewq}[$\star\star$ \iqroles{trader}]\label{iq:rc:bermudans-and-callables:4}
Pourquoi les programmes d'émission d'obligations remboursables par anticipation affectent-ils la volatilité des swaptions à longue échéance~?
\end{interviewq}

\begin{interviewq}[$\star\star\star$ \iqroles{chercheur, risque}]\label{iq:rc:bermudans-and-callables:5}
Comment valideriez-vous un pricer bermudéen~?
\end{interviewq}

\begin{interviewq}[$\star\star\star$ \iqroles{développeur}]\label{iq:rc:bermudans-and-callables:6}
Votre prix bermudéen par régression varie de deux points de base lorsque vous changez la
graine. Est-ce acceptable, et que feriez-vous~?
\end{interviewq}
